% Pillar-3 (group size) discussion + conclusion
\section{Discussion and Limitations}
\label{sec:p3-discussion}
Reframing $G$ as a preference-density dial clarifies the non-monotone effect, but
several limitations bound the claim. The held-out swings, while suggestive, are
modest against seed and difficulty variation, so we resist naming an optimal $G$;
our equivalence tests are powered to reject large differences, not to certify exact
parity. Read through the four-axis pipeline factorization of
\citet{ivison2024unpacking}, applied here to verifiable-reward stacks in the
RLVR paradigm of \citet{lambert2024tulu3}, the same-stack audit places the
algorithm axis as residual ($\eta^2{=}0.023$; paired
$\Delta_{\text{GRPO}-\text{PPO}}{=}-0.002$, permutation $p{=}0.62$) while the
$G$ axis dominates the variance budget ($\eta^2{=}0.54$ on terminal accuracy,
$0.63$ on gradient norm; Cohen's $d{=}-1.47$ for $G{=}2{\to}16$;
\texttt{experiments/results/berkeley/unpacking\_dpo\_ppo\_factorization.json}).
An independent Miller-style error-bars audit \citep{miller2024erroreval} of the
four Pillar-3 headline claims sharpens the same scoping: the
$G{=}4/G{=}32$ GU-ratio (audit H1, bootstrap CI $[4.16, 4.82]$) and the
negative retention-vs-$T$ slope (H2, CI $[-0.237, -0.038]$) are DECISIVE,
whereas the SNR-slope-vs-theory comparison (H3; CI $[+0.148, +0.583]$
contains the theoretical $+0.500$) and the native-Wu $G{=}2 \approx G{=}16$
equivalence claim (H4; $n{=}3$ paired seeds, equivalence region straddled)
are NULL
(\texttt{experiments/results/berkeley/adding\_error\_bars\_summary.json}).

% Berkeley harvest rows 01 (Dualformer, F24 L8 Yuandong Tian) and 02
% (DPO/Iterative-RPO, SP25 L2 Jason Weston); evidence:
% experiments/results/berkeley/dualformer_*.tsv,
% dpo_iterative_rpo_*.tsv + dpo_iterative_rpo_summary.json
Two further reformulations sharpen the practical reading of the $G$ axis.
First, the preference-pair view makes the DPO connection exact rather than
rhetorical: on a single winner--loser pair within a $G{=}2$ group, the GRPO
loss coincides with the small-$\beta$, no-KL, online limit of
DPO~\citep{rafailov2024direct}, and Iterative RPO's DPO-plus-NLL
objective~\citep{pang2024iterativerpo} corresponds to GRPO with replay on
winning trajectories; by analytical construction the two share the same
optimal group size at every budget ($G^\ast = 8, 16, 32, 32$ at
$T \in \{1, 4, 16, 64\}\mathrm{M}$, where the $G{=}32$ optima at high $T$
are joint-fit extrapolations---no $G{=}32$ cell and no Iterative-RPO run
was trained), and the SNR slope in $G$
($+0.366$/decade, 95\% CI $[+0.148, +0.583]$) contains the
$\sqrt{G}$-theory value of $+0.500$
(\texttt{experiments/results/berkeley/dpo\_iterative\_rpo\_summary.json}).
Second, a Dualformer-style fast/slow/auto reading of
$G$~\citep{su2024dualformer} suggests the non-monotone landscape is
exploitable: a difficulty-gated allocation rule (predicted accuracy
$\geq 0.85 \to G{=}2$, down to $G{=}32$ otherwise) attains mean
$G_{\mathrm{auto}} = 7.0$ on the 20-cell joint-fit grid --- a $56.2\%$
rollout saving versus always-$G{=}16$ and a further $\approx 42\%$ versus
the measured $G^\ast(T)$ schedule (mean $G{=}12.0$, itself a $25\%$
saving versus always-$G{=}16$) --- with 5/20 cells under-predicting by more than 5\,pp,
attributable to joint-fit residual noise rather than the allocation rule
(\texttt{experiments/results/berkeley/dualformer\_compute\_savings.tsv}).

The all-pairs contrast identity is an exact algebraic rewriting whose
observable signatures are consistent with logged per-step diagnostics
(advantage variance, ZVF, and gradient norms); direct gradient-vector validation
remains future work, and the
mechanism separating
variance-reduction from contrast-density remains only partially resolved on
existing data---the decisive controlled experiment (a token- and step-matched
$2{\times}2$ design that supplies contrast without variance and vice versa) is
specified but not yet run. Finally, the GRPO$\leftrightarrow$DPO bridge is exact
only up to a KL penalty and the projection of zero-variance groups; treating GRPO
as literally DPO ignores stabilizing structure the group baseline may carry.

\section{Conclusion}
\label{sec:p3-conclusion}
Group size in GRPO is best understood not as a variance knob to be maximized but as
a dial on how much implicit preference signal each prompt contributes. On our
benchmark its effect is non-monotone with no universal optimum, held-out
accuracy is retained almost completely across $G{=}2$ to $G{=}16$ on the
measured near-ceiling task (while the illustrative GSM8K reanalysis falls to
about $73\%$ $G{=}4$/$G{=}32$ retention at $T{=}64$M), and the signal-to-noise ratio grows
far more slowly than $\sqrt{G}$---all consistent with a contrast-density rather than
a pure-variance account, and with the exact sense in which the group-centered update
is an all-pairs preference contrast. The clean causal test, and matched-budget
multi-seed comparisons, are the natural next experiments; we release the sweeps,
per-step gradient-norm logs, and scripts to enable them. As a reproducibility
check on the released artifacts, an automated extract-and-verify pass in the
style of Paper2Agent \citep{miao2025paper2agent} recovers the analysis recipe
from the Pillar-3 headline TSVs with full field recall ($1.000$), reproduces
the published joint fit ($R^2{=}0.854$ vs.\ the published $0.796$), and
transfers to a held-out slice ($R^2{=}0.812$). The recipe also survives
deletion of a fitted field (zero failures; the slope is re-recovered from the
data alone), so the released TSVs are sufficient to reconstruct the paper's
central fit without human intervention
(\texttt{experiments/results/berkeley/paper2verifier.json}).
